% Auto-generated from book/B_open_problems.md — do not edit by hand
% Regenerated by scripts/md_to_latex.py

\chapter{Open Problems (Extended)}
\label{ch:app_b_open_problems}

Living research edges of the project. The main text (especially Chapters 3, 5--10) states each problem briefly; this appendix keeps the full statements, sandboxes, and status tracking.

\textbf{Status vocabulary:} \texttt{Open} \(\cdot\) \texttt{In progress} \(\cdot\) \texttt{Partial result} \(\cdot\) \texttt{Resolved} \(\cdot\) \texttt{Deferred}

\bigskip
\noindent\rule{\textwidth}{0.4pt}
\bigskip

\section{B.1 Status table}
\label{ch:app_b_open_problems:b-1-status-table}

\begin{center}
\small
\begin{tabular}{@{}lp{0.2\textwidth}p{0.2\textwidth}p{0.2\textwidth}p{0.2\textwidth}@{}}
\toprule
\# & Problem & Home & Status & Sandbox \\
\midrule
OP1 & Canonical quaternionic Farey structure & Ch. 3 & Open & \texttt{lib.hopf\_lattice.candidate\_adjacency} \\
OP2 & Flux topograph axioms & Ch. 5 & Open & \texttt{lib.flux\_topograph} \\
OP3 & Class number \(\leftrightarrow\) Magic Island & Ch. 6 & Open & \texttt{classify\_topograph\_type}, \texttt{class\_number\_analogue} \\
OP4 & \(Z\to\) flywheel uniqueness & Ch. 7 / 10 & Open & \texttt{map\_z\_to\_flywheel}, \texttt{stability\_landscape\_z} \\
OP5 & \(350/\pi\) first principles or falsification & Ch. 10 & Open & \texttt{lib.validation} \(\cdot\) Table T4 (Appendix D) \\
OP6 & Composition of flywheels (Gauss lift) & Ch. 8--9 & Open & \texttt{lib.composition} \(\cdot\) \texttt{lib.quaternion\_algebra} \\
\bottomrule
\end{tabular}
\end{center}


\bigskip
\noindent\rule{\textwidth}{0.4pt}
\bigskip

\section{B.2 Problem statements}
\label{ch:app_b_open_problems:b-2-problem-statements}

\subsection{OP1 --- Canonical quaternionic Farey structure}
\label{ch:app_b_open_problems:op1-canonical-quaternionic-farey-structure}

\textbf{Home:} Chapter 3

Prove uniqueness (or classify) discrete adjacency / mediant rules on the gauged Hopf lattice that reduce to classical Farey under a fixed embedding \(\mathbb{Q}\cup\{\infty\}\hookrightarrow\) lattice/base. Without a chosen primary rule, ``quaternionic Farey'' remains a family of metaphors rather than a single theory.

\textbf{Sandbox:} \texttt{candidate\_adjacency} (along-fiber phase neighbors + base angular threshold). Not claimed canonical. \textbf{Diagnostics:} \texttt{adjacency\_equivariance\_score} (Ch. 4, Exercise 4.H). \textbf{Success criteria:} equivariant rule + documented Farey reduction + comparable mediant algorithm.

\subsection{OP2 --- Flux topograph axioms}
\label{ch:app_b_open_problems:op2-flux-topograph-axioms}

\textbf{Home:} Chapter 5

Which properties of Conway topographs (separator structure, periodicity, river, face/edge values) survive for (a) quaternion norm forms and (b) lattice flux functionals? State a minimal axiom system such that Hatcher's binary case is a specialization.

\textbf{Sandbox:} \texttt{build\_flux\_topograph}, \texttt{detect\_separators}, \texttt{periodicity\_score}, \texttt{separator\_equivariance\_score}. Depends on OP1. \textbf{Success criteria:} axioms + proof or counterexample of reduction to TN Ch. 4.

\subsection{OP3 --- Class number \(\leftrightarrow\) Magic Island}
\label{ch:app_b_open_problems:op3-class-number-magic-island}

\textbf{Home:} Chapter 6

Is there a precise arithmetic invariant (class number, type number, discriminant, topological charge, ...) whose magnitude or type predicts Magic Island stability scores? Correlation is not enough; seek a structural map or a clear negative result.

\textbf{Sandbox:} \texttt{classify\_topograph\_type}, \texttt{enumerate\_reduced}, \texttt{class\_number\_analogue}, \texttt{magic\_island\_score}. Depends on OP1--OP2. \textbf{Success criteria:} invariant with out-of-sample predictive power or rigorous non-existence under stated axioms.

\subsection{OP4 --- \(Z\to\) flywheel uniqueness}
\label{ch:app_b_open_problems:op4-z-to-flywheel-uniqueness}

\textbf{Home:} Chapters 7 and 10

Up to gauge equivalence, is the map from atomic number \(Z\) to flywheel configuration unique under stated axioms? If not, classify the ambiguity and its chemical consequences.

\textbf{Sandbox:} \texttt{map\_z\_to\_flywheel[\_extended]}, \texttt{stability\_landscape\_z}. \textbf{Success criteria:} uniqueness theorem up to gauge, or explicit moduli of ambiguity.

\subsection{OP5 --- \(350/\pi\) first principles or falsification}
\label{ch:app_b_open_problems:op5-350-pi-first-principles-or-falsification}

\textbf{Home:} Chapter 10

Derive \(W_g = 350/\pi\) from lattice geometry / topological clock axioms, \textbf{or} falsify multi-domain recurrence as coincidence via pre-registered statistical tests (Table T4, Appendix D).

\textbf{Sandbox:} \texttt{table\_t4\_checklist}, \texttt{default\_hypotheses}, \texttt{run\_table\_t4\_demo}, \texttt{combine\_p\_values\_fisher}, \texttt{proximity\_to\_wg}. \textbf{Success criteria:} derivation paper, or registered negative result with full T4 compliance.

\subsection{OP6 --- Composition of flywheels (Gauss lift)}
\label{ch:app_b_open_problems:op6-composition-of-flywheels-gauss-lift}

\textbf{Home:} Chapters 8--9

Does Gauss-style composition admit a dynamical or ideal-theoretic realization on flywheels / flux classes that is associative up to gauge equivalence, compatible with Ch. 4 actions, and reduces to classical composition under restriction?

\textbf{Sandbox:} \texttt{compose\_flywheels}, \texttt{composition\_table}, \texttt{is\_associative\_up\_to\_equivalence}, \texttt{class\_group\_analogue}; algebraic side \texttt{QuaternionAlgebra}, \texttt{HurwitzOrder}, \texttt{left\_ideal\_class\_group}. \textbf{Current note:} flux-composition experiments often show \textbf{low closure and low associativity} --- document failures as progress. \textbf{Success criteria:} associative group law on classes + comparison with Hurwitz / other orders.

\bigskip
\noindent\rule{\textwidth}{0.4pt}
\bigskip

\section{B.3 How to update}
\label{ch:app_b_open_problems:b-3-how-to-update}

\begin{enumerate}
\item Change \textbf{Status} when work starts or a result lands.  
\item Put initials or issue links in an Owner column (optional).  
\item When resolved, add a one-line pointer to the theorem / section / notebook and keep the row for history.  
\item Keep \texttt{notes/open\_problems.md} in sync with this appendix (or treat this appendix as the print form of that file).
\end{enumerate}

\bigskip
\noindent\rule{\textwidth}{0.4pt}
\bigskip

\section{B.4 Dependency sketch}
\label{ch:app_b_open_problems:b-4-dependency-sketch}

\begin{lstlisting}[style=qga]
OP1 (adjacency)
  ??? OP2 (topographs)
        ??? OP3 (class number / islands)
              ??? OP4 (Z-map uniqueness)
              ??? OP6 (composition) ??? Ch. 9 ideal theory
OP5 (350/?) ?? independent empirical track, informed by geometry
\end{lstlisting}

\bigskip
\noindent\rule{\textwidth}{0.4pt}
\bigskip


